% =================================================
% Домашняя работа. Урок 2
% =================================================

\homeworktitle{Домашняя работа}{Классы. Чтение и запись чисел}
\studentline

\begin{routebox}[Как выполнять]
  В заданиях 1--2 сначала отделяйте классы справа налево. В заданиях 3--4
  следите за нулевыми и неполными классами. Задание 7 отмечено звёздочкой:
  в нём число нужно восстановить по содержимому классов.
\end{routebox}

\begin{homeworktask}[Разделите на классы]{1}
  Поставьте пробелы между классами:
  \[
    74006,\qquad 8203051,\qquad 6004007025.
  \]
  \answerspace{22mm}
\end{homeworktask}

\begin{homeworktask}[Что записано в каждом классе?]{2}
  Для числа \(408\,025\,071\) укажите:
  \begin{enumerate}
    \item число в классе миллионов;
    \item число в классе тысяч;
    \item число в классе единиц;
    \item группу цифр, которой записан класс тысяч.
  \end{enumerate}
  \answerspace{28mm}
\end{homeworktask}

\begin{homeworktask}[Прочитайте числа]{3}
  Запишите словами:
  \begin{enumerate}
    \item \(58\,014\);
    \item \(7\,006\,030\);
    \item \(320\,000\,405\);
    \item \(4\,012\,005\,008\).
  \end{enumerate}
  \answerspace{44mm}
\end{homeworktask}

\begin{homeworktask}[Запишите цифрами]{4}
  \begin{enumerate}
    \item девять миллионов сорок тысяч семь;
    \item двести пять миллионов восемнадцать;
    \item шесть миллиардов три миллиона пять тысяч девяносто;
    \item сорок миллиардов восемь.
  \end{enumerate}
  \answerspace{36mm}
\end{homeworktask}

\begin{homeworktask}[Нулевые классы]{5}
  Вставьте группы из трёх цифр:
  \begin{enumerate}
    \item \(12\,\boxed{\phantom{000}}\,305\) --- двенадцать миллионов триста пять;
    \item \(7\,042\,\boxed{\phantom{000}}\) --- семь миллионов сорок две тысячи;
    \item \(5\,\boxed{\phantom{000}}\,008\) --- пять миллиардов восемь.
  \end{enumerate}
  \answerspace{15mm}
\end{homeworktask}

\begin{homeworktask}[Найдите и исправьте ошибки]{6}
  \begin{enumerate}
    \item «тридцать миллионов пять» записали как \(30\,005\);
    \item число \(8\,004\,020\) прочитали: «восемь миллионов четыреста двадцать».
  \end{enumerate}
  Объясните каждую ошибку и дайте правильный ответ.
  \answerspace{28mm}
\end{homeworktask}

\begin{homeworktask}[Задание со звёздочкой]{7*}
  Класс миллионов содержит число \(46\), класс тысяч --- число \(5\),
  а класс единиц --- число \(8\). Запишите всё число и прочитайте его.

  \begin{hintbox}
    Каждый класс справа от первого должен занимать три позиции.
  \end{hintbox}
  \answerspace{18mm}
\end{homeworktask}

\begin{selfcheck}
  \begin{itemize}[label=\(\square\),leftmargin=1.7em,itemsep=0.25em,topsep=0.2em]
    \item Я отделяю классы справа налево.
    \item Каждый непервый класс записан тремя цифрами.
    \item Нулевые классы сохранены в записи и пропущены при чтении.
    \item После записи числа я прочитал его обратно.
  \end{itemize}
\end{selfcheck}
